\documentclass[11pt]{article}
\usepackage[margin=1in]{geometry}
\usepackage{amsmath,amssymb,amsthm}
\usepackage[T1]{fontenc}
\usepackage{lmodern}
\usepackage{hyperref}
\title{A spectral multiplicity counterexample to conjecture 00000000806}
\author{quantumtool}
\date{October 19, 2026}
\begin{document}
\maketitle

\section*{Claim and convention}
The conjecture defines $\lambda_1(\Omega)$ and $\lambda_2(\Omega)$ as the first two Dirichlet eigenvalues, and claims that the ratio
\[
R(\Omega)=\frac{\lambda_2(\Omega)-\lambda_1(\Omega)}{\lambda_1(\Omega)}
\]
attains the value $3/4$ on the union of two equal-volume disjoint disks. We use the standard spectral convention: eigenvalues are ordered nondecreasingly and repeated according to multiplicity.

\section*{Counterexample to the asserted extremizer}
Let $D$ be a disk and let $\Omega=D_1\sqcup D_2$ be the disjoint union of two congruent copies of $D$. The Dirichlet Laplacian on a disjoint union is the orthogonal direct sum of the component Laplacians. Its spectrum, counted with multiplicity, is the union of the component spectra, also counted with multiplicity.

The first Dirichlet eigenvalue $\mu_1(D)$ of a disk is simple. Since the two components are congruent, each contributes the same ground eigenvalue $\mu_1(D)$. Consequently the ground eigenvalue of $\Omega$ has multiplicity at least two, and
\[
\lambda_1(\Omega)=\lambda_2(\Omega)=\mu_1(D)>0.
\]
Therefore
\[
\gamma(\Omega)=\lambda_2(\Omega)-\lambda_1(\Omega)=0,
\qquad R(\Omega)=0,
\]
which is not $3/4$. Thus the stated union of two equal-volume disjoint disks cannot attain the claimed value under the standard eigenvalue convention, so the conjecture as written is false.

The separate assertion about connected domains does not remove this contradiction: the proposed maximizer is itself disconnected. If a different convention omitting eigenvalue multiplicity was intended, the definition of $\lambda_2$ must be changed; that is not the standard Dirichlet eigenvalue ordering used in the conjecture's displayed gap.

\section*{Lean arithmetic check}
The accompanying Lean project verifies the general arithmetic fact that two equal consecutive eigenvalues give gap zero, and that the resulting ratio is not $3/4$. The direct-sum spectral fact used to identify the two eigenvalues is established above.
\end{document}
